\section{Bộ 46}
\deda{Trong mặt phẳng tọa độ Oxy, cho điểm M(0;2) và hai đường thẳng d1:x+2y=0, d2:4x+3y=0.Viết phương trình đường tròn đi qua điểm M có tâm thuộc đường thẳng d1 và cắt d2 tại hai điểm A,B sao cho độ dài đoạn AB bằng $4\sqrt{3}$. Biết tâm đường tròn có tung độ dương.}
\hda 
Gọi $I\left ( -2t;t \right )\in d$ là tâm đường tròn $\left ( t> 0 \right )$\\
$IM^{2}=4t^{2}+\left ( 2-t \right )^{2}$\\
$d_{\left ( I,d_{2} \right )}=\left | d \right |$.\\
Gọi H là trung điểm của đoạn AB\\
Ta có: $IH^{2}+AH^{2}=IA^{2}\Leftrightarrow IH^{2}+AH^{2}=IM^{2}$\\
$\Leftrightarrow t^{2}-2t-2=0\Leftrightarrow \begin{bmatrix}
t=2\left ( tm \right ) &  & \\ 
t=-1(l) &  & 
\end{bmatrix}$\\
DO đó phương trình đường tròn là $\left ( x+4 \right )^{2}+\left ( y-2 \right )^{2}=16$


\deda{Giải hệ phương trình:
\cl{$\begin{cases}x^{3}+12y^{2}+x+2=8y^{3}+8y\\ \sqrt{x^{2}+8y^{3}}=5x-2y  \end{cases}$}}

\hda 
Điều kiện: $x^{2}+8y^{3}\geq 0$\\
Xét phương trình (1) ta có:\\
$x^{3}+12y^{2}+x+2=8y^{3}+8y\Leftrightarrow x^{3}+x+1=\left ( 2y-1 \right )^{3}+2y-1+1\Leftrightarrow x=2y-1$\\
Thế vào phương trình (2) ta được:\\
$\sqrt{x^{3}+4x^{2}+3x+1}=4x-1$\\
$\Leftrightarrow \left\{\begin{matrix}
x\geq \frac{1}{4} &  & \\ 
x^{3}-12x^{2}+11x=0 &  & 
\end{matrix}\right.\rightarrow \begin{bmatrix}
x=1\rightarrow y=1(tm)) &  & \\ 
x=11\rightarrow y=6(tm)) &  & 
\end{bmatrix}$\\
Vậy hệ phương trình có hai nghiệm (x;y) là (1;1);(11;6).\\
\deda{TÌm giá trị nhỏ nhất của biểu thức: 
$S=\frac{3}{b+c-a}+\frac{4}{a+c-b}+\frac{5}{a+b-c}$\\

Trong đó a,b,c là độ dài các cạnh của một tam giác thỏa mãn $2a+b=abc$}

\hda 
Áp dụng BĐT $\frac{1}{x}+\frac{1}{y}\geq \frac{4}{x+y}$ ta có:\\ 
$S=\left ( \frac{1}{b+c-a}+\frac{1}{a+c-b} \right )+2\left ( \frac{1}{b+c-a}+\frac{1}{a+b-c} \right )+3\left ( \frac{1}{a+c-b}+\frac{1}{a+b-c} \right )\geq \frac{2}{c}+\frac{4}{b}+\frac{6}{a}$\\
Từ giả thiết ta có: $\frac{2}{b}+\frac{1}{c}=a\rightarrow \frac{2}{c}+\frac{4}{b}+\frac{6}{a}=2\left ( a+\frac{3}{a} \right )\geq 4\sqrt{3}$\\
Vậy $S_{min}=4\sqrt{3}\Leftrightarrow a=b=c=\sqrt{3}$\\
\section{Bộ 47}
\deda{Trong mặt phẳng với hệ tọa độ Oxy cho tam giác ABC có trực tâm H(3;0) và trung điểm của BC là I(6;1). Đường thẳng AH có phương trình x+2y-3=0. GỌi D;E lần lượt là chân đường cao kẻ từ B,C của tam giác ABC. Xác định tọa độ đỉnh của tam giác ABC, biết đường thẳng DE có phương trình x-2=0 và điểm D có tùn độ dương.}
\hda  
Gọi K là trung điểm của AH. TỨ giác ADHE nội tiếp đường tròn tâm K và BCDE nội tiếp đường tròn tâm I. Suy ra $IK\perp DE$$\rightarrow IK:y-1=0$\\

Tọa độ $K\left ( 1;1 \right )\rightarrow A=\left ( -1;2 \right )$\\
Điểm $D\left ( 2;a \right )\in DE$.\\
Ta có: $KA=KD\rightarrow D=\left ( 2;3 \right )$\\
Phương trình AC:x-3y+7=0. Phương trình BC:2x-y-11=0\\
Tọa độ C(8;5) nên B=(4;-3)\\
Vậy A(-1;2);B(4;-3);C(8;5).\\

\deda{Giải hệ phương trình:\\
$\left\{\begin{matrix}
xy+2=y\sqrt{x^{2}+2}(1) &  & \\ 
y^{2}+2\left ( x+1 \right )\sqrt{x^{2}+2x+3}=2x^{2}-4x(2) &  & 
\end{matrix}\right.$}
\hda  
Vì $\sqrt{x^{2}+2}-x> \sqrt{x^{2}}-x= \left | x \right |-x\geq 0;\forall x\in \mathbb{R}$\\
Nên ta có: $\left ( 1 \right )\Leftrightarrow y=\sqrt{x^{2}+2}+x$\\

Thế vào phương trình (2) và biến đổi ta được:\\
$\left ( x+1 \right )\left [ 1+\sqrt{\left ( x+1 \right )^{2}+2} \right ]=\left ( -x \right )\left [ 1+\sqrt{\left ( -x \right )^{2}+2} \right ](\ast )$\\

Xét hàm số: $f\left ( t \right )=t\left ( 1+\sqrt{t^{2}+2} \right )$ ta có:\\
$f'\left ( t \right )=1+\sqrt{t^{2}+2}+\frac{t^{2}}{\sqrt{t^{2}+2}}> 0;\forall t\in \mathbb{R}\rightarrow x=-\frac{1}{2}$\\
Vậy hệ đã cho có nghiệm (x;y)=$\left ( -\frac{1}{2};1 \right )$\\

\deda{CHo x;y;z là các số thực dương thỏa mãn: x+y-x=-1. Tìm GTLN của biểu thức:\\
$P=\frac{x^{3}y^{3}}{\left ( x+yz \right )\left ( y+xz \right )\left ( x+xy \right )^{2}}$}
\lgi 
Ta có: $z=x+y+1\rightarrow z+xy=\left ( x+1 \right )\left ( y+1 \right )\\
x+yz=\left ( x+y \right )\left ( y+1 \right )\\
y+xz=\left ( x+y \right )\left ( x+1 \right )$\\
Do đó: $P=\frac{x^{3}y^{3}}{\left ( x+y \right )^{2}\left ( x+1 \right )^{3}\left ( y+1 \right )^{3}}\leq \frac{x^{2}y^{2}}{4\left ( x+1 \right )^{3}\left ( y+1 \right )^{3}}$\\
THeo BĐT AM-GM ta có:\\
$x+1=\frac{x}{2}+\frac{x}{2}+1\geq 3\sqrt[3]{\frac{x^{2}}{4}}\rightarrow \left ( x+1 \right )^{3}\geq \frac{27}{4}x^{2}\rightarrow 0< \frac{x^{2}}{\left ( x+1 \right )^{3}}\leq \frac{4}{27}$\\
TƯơng tự $0< \frac{y^{2}}{\left ( y+1 \right )^{3}}\leq \frac{4}{27}$\\
Vậy $P_{Min}=\frac{4}{729}\Leftrightarrow \left\{\begin{matrix}
x=y=2 &  & \\ 
z=5 &  & 
\end{matrix}\right.$\\

\section{Bộ 48}
\deda{GIải bất phương trình: $2x+5> \sqrt{2-x}\left ( \sqrt{x-1}+\sqrt{3x+4} \right )$}
\hda 
Điều kiện: $x\in \left [ 1;2 \right ]$\\
Ta có: $2x+5=(3x+4)-(x-1)=(\sqrt{3x+4}-\sqrt{x-1})(\sqrt{3x+4}+\sqrt{x-1})$\\
Đó bất phương trình tương đương:\\
$\sqrt{3x+4}-\sqrt{x-1}> \sqrt{2-x}$ (Vì $\sqrt{3x+4}+\sqrt{x-1}> 0$)\\
CHuyển vế bình phương giải ta được nghiệm của BPT là $S=\left [ 1;2 \right ]$\\

\deda{Trong mặt phẳng tọa độ Oxy cho hình thang ABCD có $\angle BAD=\angle ADC=90^{\circ}$, AB=AD=2; DC=4, đỉnh C nằm trên đường thẳng d:2x-y+2=0. Điểm M nằm trên cạnh AD sao cho AM=2MD và đường thẳng BM có phương trình là 3x-2y+2=0. TÌm tọa độ C.}
\hda  
Ta có: $C=\left ( t;2+3t \right )\in d\rightarrow d_{\left ( C;BM \right )}=\frac{\left | 2+3t \right |}{\sqrt{13}}$\\
Theo giả thiết ta suy ra $MD=\frac{1}{3}AD=\frac{2}{3};AM=\frac{4}{3}$\\
Tam giác ABM vuông tại A nên $S_{\Delta ABM}=\frac{1}{2}AM.AB=\frac{4}{3}$\\
Tam giác CDM vuông tại D nên $S_{\Delta CDM}=\frac{1}{2}MD.DC=\frac{4}{3}$\\
Và $S_{ABCD}=\frac{1}{2}AD.(AB+CD)=6$\\
Do đó: $S_{BMC}=\frac{10}{3}$\\
Lại có: $S_{BMC}=\frac{1}{2}BM.d_{\left ( C;BM \right )}\rightarrow d_{\left ( C;BM \right )}=\frac{10}{\sqrt{13}}$\\
.....\\
\vspace{5mm}
 

 

\hspace{12cm} Đáp số: $C\left ( -4;-10 \right )$ hoặc $C=(\frac{8}{3};10)$\\
\deda{Cho a;b;c là các số thực dương thỏa mãn: $3\left ( a^{2}+b^{2}+c^{2} \right )=1$. TÌm GTNN của biểu thức:\\
$Q=\sqrt{a^{2}+b^{2}+\frac{1}{b^{2}}+\frac{1}{c^{2}}}+\sqrt{b^{2}+c^{2}+\frac{1}{c^{2}}+\frac{1}{a^{2}}}+\sqrt{c^{2}+a^{2}+\frac{1}{a^{2}}+\frac{1}{b^{2}}}$}
\hda  
Từ giả thiết suy ra $a+b+c\leq 1$\\
THeo BĐT Cau chy-Schwarz ta có:\\
$Q\sqrt{2}=\sum \sqrt{2\left ( a^{2}+b^{2} \right )+2\left ( \frac{1}{b^{2}}+\frac{1}{c^{2}} \right )}\geq \sum \sqrt{\left ( a+b \right )^{2}+\left ( \frac{1}{b}+\frac{1}{c} \right )^{2}}$\\

Xét các vecto $\vec{x}=\left ( a+b;\frac{1}{b}+\frac{1}{c} \right )
\vec{y}=\left (b+c;\frac{1}{c}+\frac{1}{a}\right)
\vec{z}=\left(c+a;\frac{1}{a}+\frac{1}{b}\right)$\\
Ta có: $\left|\vec{x}\right|+\left | \vec{y} \right |+\left | \vec{z} \right |\geq \left | \vec{x}+\vec{y}+\vec{z} \right |$\\
Do đó: $Q\sqrt{2}\geq 2\sqrt{\left ( a+b+c \right )+\frac{81}{\left ( a+b+c \right )^{2}}}$\\
Đặt $t=a+b+c;0< t\leq 1$.\\
Xét hàm số: $f\left ( t \right )=t+\frac{81}{t};t\in \left(0;1\right]$\\
.....\\
\vspace{5mm}
\hspace{12cm}$Q_{Min}=2\sqrt{41}\Leftrightarrow a=b=c=\frac{1}{3}$\\

\section{Bộ 49}
\deda{Giải hệ phương trình:
\begin{center}$\left\{\begin{matrix}
\left ( \sqrt{y}+1 \right )^{2}+\frac{y^{2}}{x}=y^{2}+2\sqrt{x-2} \hspace{2mm}(1)&  & \\ 
x+\frac{x-1}{y}+\frac{y}{x}=y^{2}+y \hspace{2mm}(2) &  & 
\end{matrix}\right.$\end{center}}
\hda 
Ta có: $(2)\Leftrightarrow \left ( x-y^{2} \right )\left ( xy+x-1 \right )=0\Leftrightarrow x=y^{2}$ \hspace{2mm}(xy+x-1>0)\\
Do đó: $(1)\Leftrightarrow \left ( \sqrt{y}+1 \right )^{2}=\left ( \sqrt{y^{2}-2}+1 \right )^{2}\Leftrightarrow y=2\rightarrow x=4$\\

\deda{Trong mặt phẳng tọa độ Oxy, họi H(3;-2),I(8;11),K(4;-1) lần lượt là trực tâm, tâm đường tròn ngoại tiếp, chân đường cao vẽ từ A của tam giác ABC. Tìm tọa độ A;B;C.}
\hda 
$\vec{HK}=(1;1)\Rightarrow AK:x-y-5=0;BC:x+y-3=0$\\

GỌi M là trung điểm của BC suy ra $IM\perp BC\rightarrow IM:x-y+3=0\rightarrow M=(0;3)$\\
$\vec{HA}=2\vec{MI}=(16;16)\rightarrow A=(19;14)$\\

Điểm B(b;3-b)$\in BC\Rightarrow C=(-b;b+3)\rightarrow \vec{BH}=(3-b;b-5);\vec{CA}=(19+b;11-b)$\\
Ta có: $\vec{BH}\perp \vec{AC}\rightarrow \vec{BH}.\vec{CA}=0\Leftrightarrow b=\pm 1$\\
$\begin{bmatrix}
B=(1;2);C=(-1;4) &  & \\ 
B=(-1;4);C=(1;2 &  & 
\end{bmatrix}$\\

\deda{Cho hai số thực x;y thỏa mãn điều kiện $x^{4}+16y^{4}+2\left ( 2xy-5 \right )^{2}=41$. TÌm GTNN-GTLN của biểu thức: $P=xy-\frac{3}{x^{2}+4xy^{2}+3}$}
\hda 
Ta có: $x^{4}+16y^{4}+2\left ( 2xy-5 \right )^{2}=41\Leftrightarrow \left ( x^{2}+4y^{2} \right )^{2}+9=40xy$\\
Đặt $t=x^{2}+4y^{2}\rightarrow t^{2}+9=40xy\leq 10t\Rightarrow 1\leq t\leq 9$\\
Khi đó: $P=\frac{t^{2}+9}{40}-\frac{3}{3+t}$\\
.....\\
Vậy: $P_{Max}=2\Leftrightarrow x=\frac{3}{\sqrt{2}};y=\frac{3}{2\sqrt{2}}\\
P_{Min}=-\frac{1}{2}\Leftrightarrow x=\frac{1}{\sqrt{2}};y=\frac{1}{2\sqrt{2}}$\\

\section{Bộ 50}
\deda{Giải hệ phương trình\\
\begin{center}$\left\{\begin{matrix}
x^{2}y+x^{2}+1=2x\sqrt{x^{2}y+2}\hspace{5mm} (1)&  & \\ 
y^{3}\left ( x^{6}-1 \right )+3y\left ( x^{2}-2 \right )+3y^{2}+4=0 \hspace{5mm}(2)&  & 
\end{matrix}\right.$\end{center}}
\hda 
*ĐKXĐ: $x^{2}y\geq -2$\\

Ta có: $(2)\Leftrightarrow x^{6}y^{3}+3x^{2}y=y^{3}-3y^{2}+3y-1+3\left ( y-1 \right )\\
\Leftrightarrow \left ( x^{2}y \right )^{3}+3x^{2}y+3x^{2}y=\left ( y-1 \right )^{3}+3\left ( y-1 \right )$\\

Xét hàm số: $f\left ( t \right )=t^{3}+3t$ $\rightarrow$ hàm số đồng biến $\rightarrow$ $x^{2}y=y-1\hspace{5mm}(y\geq -1)$. Thế vào phương trình (1 ta được:\\
$x^{2}y+x^{2}+1=2x\sqrt{y+1}\Leftrightarrow \left ( x\sqrt{y-1}-1 \right )^{2}=0\Leftrightarrow x\sqrt{y+1}=1$\\
Do đó hệ đã cho tương đương với: $\left\{\begin{matrix}
x\sqrt{y-1}=1 &  & \\ 
x^{2}y=y-1 &  & 
\end{matrix}\right.$\\
Bình phương phương trình trên và giải cho ta nghiệm $(\frac{1+\sqrt{5}}{2};\frac{-1+\sqrt{5}}{2});(\frac{1-\sqrt{5}}{2};\frac{1+\sqrt{5}}{2})$\\
\deda{Trong mặt phẳng hệ tọa độ Oxy, cho điểm E(3;4) đường thẳng $d:2+y-1=0$ và đường tròn $(C):x^{2}+y^{2}+4x-2y-4=0$. Gọi M là điểm thuộc đường thẳng d và nằm ngoài đường tròn (C). Từ M kẻ tiếp tuyến MA, MB đến đường tròn (C) (A, B là các tiếp điểm). Gọi (E) là đường tròn tâm E và tiếp xúc cới đường thẳng AB. Tìm tọa độ điểm M sao cho đường tròn (E) có chu vi lớn nhất.}
\hda 
ĐƯờng tròn (C) có tâm I(-2;1) và bán kính E=3.\\
Điểm M(a;1-a) thuộc d và  nằm ngoài (C) $\rightarrow$ IM>R$\rightarrow$$IM^{2}>9\Rightarrow(a+2)^{2}+(-a)^{2}>9\rightarrow2a^{2}+4a-5>0$.\\
Ta có: $MA^{2}=MB^{2}=IM^{2}-IA^{2}=\left ( a+2 \right )^{2}+\left ( -a \right 
)^{2}-9=2a^{2}+4a-5$\\
Do A;B thuộc đường tròn (C) nên tọa độ A;B thỏa mãn phương trình: $x^{2}+y^{2}+4x-2y-4=0$\\
Suy ra: $(a+2)x-ay+3a-5=0$\hspace{5mm}(3)\\
(3) là phương trình đường thẳng AB\\
Đường tròn (E) có bán kính là $R'=d_{E;\Delta }$\\
DO đó chu vi (E) lớn nhất $\leftrightarrow$R' lớn nhất $\leftrightarrow$ $R'=d_{E;\Delta }$ lớn nhất\\
Đường thẳng $\Delta$ luôn đi qua điểm K($\frac{5}{2};\frac{11}{2}$)\\
GỌi H là hình chiếu của E lên $\Delta$. Suy ra $\$d_{E;\Delta }=EH\leq EK=\frac{\sqrt{10}}{2}$\\
Dấu "=" xảy ra khi $\Delta \perp EK$\\
....
Vậy M=(-3;4)\\
\deda{Cho $x;y;z$ là các số thực thỏa mãn: $-1-2\sqrt{2}< x< -1+2\sqrt{2};y> 0;z> 0;x+y+z=-1$. Tìm GTNN của biểu thức:
$$P=\frac{1}{\left ( x+y \right )^{2}}+\frac{1}{\left ( x+z \right )^{2}}+\frac{1}{8-\left ( y+z \right )^{2}}$$}
\hda 
Ta có: $P=\frac{1}{\left ( -1-z \right )^{2}}+\frac{1}{\left ( -1-y \right )^{2}}+\frac{1}{8-\left ( -1-x \right )^{2}}$\\
Bằng phương pháp biến đổi tương đương ta dễ dàng chứng minh được:\\
$\frac{1}{\left ( 1+y \right )^{2}}+\frac{1}{\left ( 1+z \right )^{2}}\geq \frac{1}{1+yz}$\\
Ta lại có: $yz\leq \frac{\left ( y+z \right )^{2}}{4}=\frac{\left ( 1+x \right )^{2}}{4}\rightarrow \frac{1}{\left ( 1+y \right )^{2}}+\frac{1}{\left ( 1+z \right )^{2}}\geq \frac{4}{4+\left ( 1+x \right )^{2}}\\
\Rightarrow P\geq \frac{4}{4+\left ( 1+x \right )^{2}}+\frac{1}{8-\left ( x+1 \right )^{2}}$
Do $-1-2\sqrt{2}< x< -1+2\sqrt{2}\rightarrow 0\leq \left ( x+1 \right )^{2}< 8$\\
Đến đây ta xét hàm số và tìm được $P_{Min}=\frac{3}{4}\Leftrightarrow x=-3;y=z=1$\\

\section{Bộ 51}
\deda{Giải hệ phương trình:\\
\begin{center}$\left\{\begin{matrix}
\sqrt{x-1}+\sqrt{y-1}=2 &  & \\ 
\sqrt{x+2}+\sqrt{y+2}=4 &  & 
\end{matrix}\right.$\end{center}}
\hda 
Hệ đã cho tương đương với: $\left\{\begin{matrix}
x-1+y-1+2\sqrt{\left ( x-1 \right )\left ( y-1 \right )}=4&  & \\ 
x+2+y+2+2\sqrt{\left ( x+2 \right )\left ( y+2 \right )}=16 &  & 
\end{matrix}\right.\\
\Leftrightarrow \left\{\begin{matrix}
(x-1)+(y-1)+2\sqrt{\left ( x-1 \right )\left ( y-1 \right )}=4\hspace{8mm}(3) &  & \\ 
\left ( x-1 \right )+(y-1)+2\sqrt{(x-1)(y-1)+3(x-1+y-1)+9}=10\hspace{8mm}(4) &  & 
\end{matrix}\right.$\\
Đặt $t=\sqrt{(x-1)(y-1)}\geq 0\rightarrow x-1+y-1=4-2t$\\
Thế vào (4) ta được $\sqrt{t^{2}-6t+21}=t+3\Leftrightarrow t=1\Rightarrow \left\{\begin{matrix}
x=2 &  & \\ 
y=2 &  & 
\end{matrix}\right.$\\
\deda{Trong mặt phẳng tọa độ Oxy cho tam giác ABC có AB=3AC. ĐƯờng phân giác trong của góc BAC có phương trình x-y=0. Đường cao BH có phương trình: 3x+y-16=0. Hãy xác định tọa độ A,B,C biết AB đi qua điểm M(4;10)}
\hda 
Gọi M' là điểm đối xứng với M qua đường phân giác trong góc BAC. SUy ra M'=(10;4) và thuộc đường AC.\\
Vì đường cao qua đỉnh B có phương trình 3x+y-16=0 nên AC: x-3y+2=0\\
Suy ra: A=(1;1)$\rightarrow$AB:3x-y-2$\rightarrow$B=(3;7)\\
Mặt khác AB=3AC $\Rightarrow$ C=(3;$\frac{5}{3}$)\\
Vậy A(1;1);B(3;7);C(3;$\frac{5}{3}$)
\deda{Xét 3 số thực x;y;z thỏa mãn $x^{3}+y^{3}+z^{3}-3xyz=1$. TÌm GTNN của biểu thức: $P=x^{2}+y^{2}+z^{2}$}
\hda 
Ta có: $1=x^{3}+y^{3}+z^{3}-3xyz=\left ( x+y+z \right )\left ( x^{2}+y^{2}+z^{2}-xy-yz-xz \right )\hspace{5mm}(1)\\
\Rightarrow x^{2}+y^{2}+z^{2}-xy-yz-xz> 0\rightarrow x+y+z> 0$\\
Từ (1) ta được $P=\frac{t^{2}}{3}+\frac{2}{3t}$ với t=x+y+z>0\\
Vậy $P_{Min}=1\Leftrightarrow \left ( x;y;z \right )=(1;0;0)$ và các hoán vị.
\section{Bộ 52}
\deda{Giải phương trình:
$2x^{3}+9x^{2}-6x(1+2\sqrt{6x-1})+2\sqrt{6x-1}+8=0$}
\hda 
ĐK: $x\geq \frac{1}{6}$\hspace{7mm}(1)\\
Ta có: $(1)\Rightarrow2x^{3}+9x^{2}-6x+8=2(6x-1)\sqrt{6x-1}$\\
Đặt $y=\sqrt{6x-1}$.Ta có hệ phương trình:\\
$\left\{\begin{matrix}
2x^{3}+9x^{2}-6x+8=2y^{3} &  & \\ 
18x-3=3y^{2} &  & 
\end{matrix}\right.\\
\Rightarrow 2x^{3}+9x^{2}+12x+5=2y^{3}+3y^{2}\\
\Leftrightarrow 2(x+1)^{3}+3(x+1)^{2}=2y^{3}+3y^{2}$\\
Xét hàm số: $f(t)=2t^{3}+2t^{2};t\geq 0$. Hàm số đồng biến trên $\left[0;\infty\right)$\\
Do đó: x+1=y.\\
Vậy phương trình đã cho có 2 nghiệm $x=2+\sqrt{2};x=2-\sqrt{2}$\\
\deda{Trong mặt phẳng hệ tọa độ Oxy cho tam giác ABC vuông tại A, có B(-2;1) và C(8;1). Đường tròn nội tiếp tam giác ABC có bán kính $r=3\sqrt{5}-5$. Tìm tọa độ tâm đường tròn nội tiếp tam giác ABC biết $y_{I}>0$.}
\hda 
Gọi p là nửa chu vi tam giác ABC.\\
Ta có BC=10. Gọi M,N lần lượt là tiếp điểm trên AB,AC$\rightarrow p=BC+AM=3\sqrt{5}+5$. Ta có $S=pr=20$\\
Đặt AH=h. Ta có $S_{ABC}=\frac{1}{2}BC.h=20\rightarrow h=4$\\
Do $r=3\sqrt{5}-5$ nên I nằm trên đường thẳng song song với BC cách BC 1 đoạn đúng bằng r. Mà $y_{I}>0$ nên I nằm trên đường $y=3\sqrt{5}-4$ và A nằm trên đường y=5.\\ 
Gọi J là trung điểm BC suy ra JA=IB=JC $\Rightarrow A(0;5) or A(6;5)$\\

\deda{Cho 2 số thực dương tùy ý a;b;c. Tìm GTNN của biểu thức:

\begin{center}
$P=\frac{\sqrt{a^{3}c}}{2\sqrt{b^{3}a}+3bc}+\frac{\sqrt{b^{3}a}}{2\sqrt{c^{3}b}+3ca}+\frac{\sqrt{c^{3}b}}{2\sqrt{a^{3}c}+3ab}$
\end{center}}
\hda 
Ta có: $\frac{\sqrt{a^{3}c}}{2\sqrt{b^{3}a}+3bc}=\frac{a\sqrt{ac}}{b(2\sqrt{ba}+3c)}=\frac{(\sqrt{\frac{a}{b}})^{2}}{2\sqrt{\frac{b}{c}}+3\sqrt{\frac{c}{a}}}$\\
Thiết lập các biểu thức tương tự.\\
Đổi biến: $(\sqrt{\frac{a}{b}};\sqrt{\frac{b}{c}};\sqrt{\frac{c}{a}})\rightarrow (x;y;z)\rightarrow xyz=1$\\
Ta được: $P=\frac{x^{2}}{2y+3z}+\frac{y^{2}}{2z+3x}+\frac{z^{2}}{2x+3y}\geq \frac{1}{5}\left ( x+y+z \right )\geq \frac{3}{5}\sqrt[3]{xyz}=\frac{3}{5}$\\
Vậy $P_{Min}=\frac{3}{5}\Leftrightarrow x=y=z=1$ hay $a=b=c$\\
\section{Bộ 53}
\deda{Giải hệ phương trình trên tập số thực:
\begin{center}
$\left\{\begin{matrix}
2\sqrt{x^{2}+5}=2\sqrt{2y}+x^{2}\hspace{7mm}(1) &  & \\ 
x+3\sqrt{xy+x-y^{2}-y}=5y+4\hspace{7mm}(2) &  & 
\end{matrix}\right.$
\end{center}}
\hda 
ĐKXĐ:\\ 
Ta có: $(2)\Leftrightarrow (x-2y-1)+3\left ( \sqrt{xy+x-y^{2}-y}-y-1 \right )=0\\
\Leftrightarrow (x-2y-1)\left [ \frac{3(y+1)}{\sqrt{xy+x-y^{2}-y}+y+1} \right ]=0\\
\Leftrightarrow x-2y-1=0$\hspace{5mm}(Vì theo ĐKXĐ thì trong ngoặc >0)\\
Thế vào (1) ta được:\\

$2\sqrt{x^{2}+5}=2\sqrt{x-1}+x^{2}\\
\Leftrightarrow (x-2)\left [ -\frac{2(x+2)}{\sqrt{x^{2}+5}+3}+\frac{2}{\sqrt{x-1}+1}+x+2 \right ]=0$\hspace{5mm}(3)\\
Với $\forall x\geq 1$ thì:\\
 $-\frac{2(x+2)}{\sqrt{x^{2}+5}+3}+\frac{2}{\sqrt{x-1}+1}+x+2\\
=\frac{2}{\sqrt{x-1}+1}+(x+2)(1-\frac{2}{\sqrt{x^{2}+5}+3})> 0$\\
Vậy hệ phương trình đã cho có 1 nghiệm (x;y)=(2;$\frac{1}{2}$)\\
\deda{Trong mặt phẳng tọa độ Oxy cho tứ giác ABCD nội tiếp đường tròn AC. Điểm M(3;-1) là trung điểm của BD, C(4;-2). Điểm N(-1;-3) nằm trên đường thẳng đi qua B và vuông góc với AD. Đường thẳng AD đi qua P(1;3). Tìm tọa độ A;B;D}
\hda 
Giả sử D(a;b). Vì M là trung điểm của BD nên B=(6-a;-2-b).\\
Ta có: $\angle ADC=90^{\circ}\rightarrow AD\perp DC\rightarrow BN\parallel CD$\\
Vì $\vec{NB};\vec{CD}$ cùng phương nên ta được b=a-6\hspace{5mm}(1)\\
Lại có: $\vec{PD}\perp \vec{CD}\rightarrow (a-1)(a-4)+(b+2)(b-3)=0$\hspace{5mm}(2)\\
Thế (1) vào (2) ta suy ra a=5 hoặc a=4
....\\
.\hspace{10cm}\underline{Đáp số:}A(2;2);D(5;-1);B(1;-1)\\
\deda{Cho x là số thực thuộc đoạn $\left [ -1;\frac{5}{4} \right ]$. TÌm GTLN, GTNN của biểu thức:
\begin{center}
$P=\frac{\sqrt{5-4x}-\sqrt{1+x}}{\sqrt{5-4x}+2\sqrt{1+x}+6}$
\end{center}}
\hda 
Đặt $\left\{\begin{matrix}
\sqrt{5-4x}=a &  & \\ 
\sqrt{1+x}=b &  & 
\end{matrix}\right.\Rightarrow a^{2}+4b^{2}=9$.
Do đó tồn tại $\alpha \in \left [ 0;\frac{\pi }{2} \right ]\Rightarrow \left\{\begin{matrix}
a=3\sin\alpha  &  & \\ 
2b=3\cos\alpha  &  & 
\end{matrix}\right.$\\
Khi đó $P$ trở thành:
$P=\frac{2\sin\alpha -\cos\alpha }{2\sin\alpha +2\cos\alpha +4}$\\
Đến đây ta sẽ xét hàm số.\\
.\hspace{10cm}\underline{Đáp số:}$\left\{\begin{matrix}
P_{Min}=\frac{-1}{6}\Leftrightarrow x=\frac{5}{4} &  & \\ 
P_{Max}=\frac{1}{3}\Leftrightarrow  &  & 
\end{matrix}\right.$\\

\section{Bộ 54}
\deda{Giải hệ phương trình:
\begin{center}
$\left\{\begin{matrix}
x^{2}+6y-4=\sqrt{2(1-y)(x^{3}+1)}\hspace{5mm}(1) &  & \\ 
(3-x)\sqrt{2-x}-2y\sqrt{2y-1}=0\hspace{5mm}(2) &  & 
\end{matrix}\right.$
\end{center}}
\hda 
ĐKXĐ:\\
Ta có:\\
$(2)\Leftrightarrow (1+2-x)\sqrt{2-x}=(1=2y-1)\sqrt{2y-1}$\\
Xét hàm số: $f(t)=(1+t^{2})t=t^{3}+t$\\
Hàm số đồng biến nên suy ra: 2y=3-x\\

Thế cà (2) biến đổi thành: $3(x^{2}-x+1)-2(x^{2}-1)=\sqrt{(x^{2}-1)(x^{2}-x+1)}\Leftrightarrow x=2$\\
\deda{Trong mặt phẳng tọa độ Oxy cho hình chữ nhật ABCD có diện tích bằng 12, I($\frac{9}{2}$;$\frac{3}{2}$) là tâm hình chữ nhật và M(3;0) là trung điểm của AD. Tìm tọa độ các đỉnh của hình chữ nhật biết $y_{D}$<0}
\hda
\vspace{5cm}

\deda{Cho các số thực dương a;b;c thỏa mãn: $a^{2}+b^{2}+c^{2}=1$. TÌm GTNN của biểu thức:$P=\frac{1}{\sqrt{a^{2}+ab}}+\frac{1}{\sqrt{b^{2}+ab}}+\frac{2\sqrt{3}}{1+c}$}
\hda 
Theo BĐT Cauchy ta có: 
$\frac{1}{\sqrt{a^{2}+ab}}+\frac{1}{\sqrt{b^{2}+ab}}\geq \frac{2}{\sqrt[4]{(a^{2}+ab)(b^{2}+ab)}}\geq \frac{2}{\sqrt{\frac{a^{2}+ab+b^{2}+ab}{2}}}\geq \frac{2}{\sqrt{a^{2}+b^{2}}}$\\
Do đó: $P\geq \frac{2}{\sqrt{1-c^{2}}}+\frac{2\sqrt{3}}{c+1}$\\
Vậy $P_{Min}=\frac{8\sqrt{3}}{3}\Leftrightarrow a=b=\frac{\sqrt{6}}{4};c=\frac{1}{2}$\\
\section{Bộ 55}
\deda{Giải bất phương trình:$\sqrt{4x+1}+\sqrt{6x+4}\geq 2x^{2}-2x+3$}
\hda 
$bpt\Leftrightarrow \sqrt{4x+1}-(x+1)+\sqrt{6x+4}-(x+2)\geq 2(x^{2}-2x)\\
\Leftrightarrow (x^{2}-2x)(2+\frac{1}{\sqrt{4x+1}+x+1}+\frac{1}{\sqrt{6x+4}+x+2})\geq 2(x^{2}-2x)$
\deda{Trong mặt phẳng Oxy cho tam giác ABC có A(1;5), đường phân giác trong góc A có phương trình x-1=0, tâm đường tròn ngoại tiếp tam giác ABC là I($\frac{-3}{2}$;0) và điểm M(10;2) thuộc đường thẳng BC. Tìm tọa độ B;C}
\hda
\vspace{5cm}

\deda{Cho a;b;c là 3 số dương thỏa mãn $s^{2}+b^{2}+c^{2}=1$. TÌm GTNN của biểu thức: $P=\frac{1}{a^{4}+a^{2}b^{2}}+\frac{1}{b^{4}+a^{2}b^{2}}+\frac{32}{(1+c)^{3}}$}
\hda 
Theo BĐT Cauchy ta có: $\frac{1}{a^{4}+a^{2}b^{2}}+\frac{1}{b^{4}+a^{2}b^{2}}\geq \frac{4}{\left ( a^{2}+b^{2} \right )^{2}}= \frac{4}{\left ( 1-c^{2} \right )^{2}}\\
\Rightarrow P\geq \frac{4}{\left ( 1-c^{2} \right )^{2}}+\frac{32}{\left ( 1+c \right )^{3}}$\\
....
Vậy ${Min}=\frac{448}{27}\Leftrightarrow a=b=\frac{\sqrt{6}}{4};c=\frac{1}{c}$